\subsection{二重覆盖中的消失圈与匹配圈}\label{sec:18c}
取 $\pi$ 的消失路径 $\gamma:[0,1]\to S$，从 $\gamma(0)=*$ 出发，并始终落在较小的圆盘 $\mathcal W\subset S$ 内。由于 $\Psi:\Efamily\to\mathcal W$ 与 $\pi$ 有相同的临界值，也可将 $\gamma$ 视为该 Lefschetz 纤维化的消失路径。因此，有相应的消失圈 $\widetilde V_\gamma$，它是 $\widetilde E=\Efamily_*$ 中的 Lagrangian 球面。另一方面，已在\secref{18a}中看到，$\gamma$ 的原像经适当参数化成为 $\widetilde\gamma:[-1,1]\to\widetilde S$，是 $\widetilde\pi$ 的匹配路径，其相应匹配圈为 $\widetilde\Delta_\gamma$。\znote{272}{这里用了 §16g 的匹配圈构造：匹配路径 $\widetilde\gamma$ 的两个半段各有一个 Lefschetz 顶针，二者的公共边界恰是消失圈（此处即 $M$ 中 $\gamma$ 的消失圈），沿该圈把两个顶针粘合，所得球面就是匹配圈。于是同一族消失路径有两种读法：把 $\gamma$ 看成 $\Psi$ 的消失路径，得到 $\widetilde E$ 中的球面 $\widetilde V_\gamma$；把 $\widetilde\gamma$ 看成 $\widetilde\pi$ 的匹配路径，得到 $\widetilde\Delta_\gamma$。引理 18.2 断言二者只差带标架同痕，这正是全节把 (B) 型膜同时写成顶针与匹配圈的合法性所在。[来源：原书 §16g、§18c（印刷页 272）]}

\StatementNumber{2}\begin{lemma}\label{st:18.2}
作为 $\widetilde E$ 中的带标架 Lagrangian 球面，$\widetilde\Delta_\gamma$ 与 $\widetilde V_\gamma$ 同痕。
\end{lemma}
\begin{proof}
其中大部分内容可以从双纤维化的一般理论，具体而言从\stref{引理}{16.15}推出。考察 $\operatorname{Critv}(\varpi)\subset\Sfamily$ 的结构，见\eref{18.6}，就能看出，与 $\gamma$ 相关的匹配路径与 $\widetilde\gamma$ 同痕。这一论证唯一的问题是，\stref{引理}{16.15}没有刻画其中出现的同痕 $I$。当然，在当前情形中，这个同痕有些多余，因为匹配圈可以用朴素方式构造；不过，仍值得核查各个要素确实能够正确结合。与往常一样，先在唯一临界点 $x\in E_z$，$z=\gamma(1)$，附近作局部变形，使 Kähler 形式 $\omega_E$ 在局部全纯 Morse 坐标下成为标准形式。还可假设，当 $t$ 接近 $1$ 时，$\gamma(t)=z+(1-t)$。这两种改变都允许，因为它们不影响所考虑的两个带标架 Lagrangian 球面的同痕类。取 $s$ 充分接近 $1$，考虑短消失路径 $\gamma|[s,1]$，及其在 $\Sfamily_{\gamma(s)}\to S$ 下的原像；该原像是 $\varpi_{\gamma(s)}$ 的匹配路径。由假设，这一原像实际上就是连接 $z-\gamma(s)$ 的两个平方根的直线段。

取\eref{18.7}中的局部坐标 $(y,x)$，使 $\omega_{\Efamily}=d\kappa+\omega_E+\Psi^*\omega_{\mathcal W}$，其中第二项现在是 $x$ 坐标中的标准形式。这并不完全是 $\CC^{n+2}$ 上的标准形式，但由 $y\in i\R$、$x\in\R^{n+1}$ 定义的子空间，仍是 $\omega_{\Efamily}$-Lagrangian 的。根据前面对 Lefschetz 顶针的刻画，容易看出，该子空间的适当一部分必为 $\gamma|[s,1]$ 的 Lefschetz 顶针。因此，消失圈为
\begin{equation}\label{eq:18.13}
\widetilde V_{\gamma|[s,1]}=\{y\in i\R,\ x\in\R^{n+1}:|y|^2+|x|^2=\gamma(s)-z\},
\end{equation}
\SourcePage{280}
并带有显然的标架。到目前为止，$\gamma|[s,1]$ 都被视为 $\Psi$ 的消失路径。若暂时再将其看作 $\pi$ 的消失路径，相应的 Lefschetz 顶针当然是 $\{x\in\R^{n+1}:|x|^2\le\gamma(s)-z\}$；因此，这个顶针在 $\beta_{\Efamily}|\Efamily_{\gamma(s)}$ 下的原像，按定义就是 $\widetilde\Delta_{\gamma|[s,1]}$，与\eref{18.13}相同。再像\stref{引理}{16.15}中那样，使用适当的平行移动映射，将这一事实传回整条路径 $\gamma$。
\end{proof}

为便于后面引用，再考察另外两种几何情形，其中仍涉及上面定义的各类 Lagrangian 子流形。假设给定 $\pi$ 的三条消失路径 $\gamma_k,\gamma_{k+1},\delta_k$，如\sr{fig:16.2}{图~16.2}所示；当然，此时不需要整个特选基，不过为方便起见仍沿用该记号。以 $\widetilde\Delta_k,\widetilde\Delta_{k+1},\widetilde\Pi_k$ 表示 $\widetilde E$ 中相应的 (B) 型 Lagrangian 子流形。
\StatementNumber{3}\begin{lemma}\label{st:18.3}
作为 $\widetilde E$ 中的带标架 Lagrangian 球面，$\widetilde\Pi_k$ 与 $\tau_{\widetilde\Delta_k}(\widetilde\Delta_{k+1})$ 同痕。
\end{lemma}
\begin{proof}
不失一般性，可假设这些路径包含于 $\mathcal W$。由\stref{引理}{18.2}，将 $\gamma_k$ 看成 $\Psi$ 的消失路径后，$\widetilde\Delta_k$ 就是其消失圈。因此，相应的 Dehn 扭转是 $\Psi$ 沿路径 $\lambda$ 的单值变换，其中 $\lambda$ 沿该消失路径绕去再绕回，即
\[
\tau_{\widetilde\Delta_k}(\widetilde\Delta_{k+1})\simeq h_\lambda(\widetilde\Delta_{k+1}).
\]
按定义，$\delta_k$ 由取 $\gamma_{k+1}$、再将其起点沿 $\lambda$ 移动得到。这意味着存在 $\pi$ 的消失路径的光滑族 $\gamma^s$，满足 $\gamma^s(0)=\lambda(s)$、$\gamma^s(1)=\gamma_{k+1}(1)$，从 $\gamma^0=\gamma_{k+1}$ 开始，到 $\gamma^1=\delta_k$ 结束。它们给出带标架 Lagrangian 球面的光滑族 $\widetilde\Delta^s\subset\Efamily_{\lambda(s)}$，在 $\widetilde\Delta^0=\widetilde\Delta_{k+1}$ 与 $\widetilde\Delta^1=\widetilde\Pi_k$ 之间插值。考察这些球面沿 $\lambda|[s,1]$ 平行移动后的像，即得到 $h_\lambda(\widetilde\Delta_{k+1})$ 与 $\widetilde\Pi_k$ 之间的同痕。
\end{proof}

最后，设 $\mu$ 为 $\pi$ 的匹配路径，$\Sigma_\mu\subset E$ 为其匹配圈。由假设，负 Liouville 流将每个 Lagrangian 子流形移入 $E^{\mathrm{shrink}}$；又因为 $\Sigma_\mu$ 只在相差同痕的意义下定义，故不妨假设它已经位于 $E^{\mathrm{shrink}}$ 内。这样，其原像 $\widetilde\Sigma_\mu=\widetilde\Sigma_{\mu,+}\cup\widetilde\Sigma_{\mu,-}$ 就是 (U) 型 Lagrangian 子流形，由 $\widetilde E$ 中两个不交的带标架 Lagrangian 球面组成。选择从 $*$ 出发、分别终止于 $\mu$ 两端的消失路径 $\gamma_k,\gamma_{k+1}$，如\sr{fig:18.1}{图~18.1}所示。关键要求是，这三条路径围成的三角形内部不含 $\pi$ 的临界值，从而复合路径 $\gamma'_{k+1}=\mu\#\gamma_k$ 是与 $\gamma_{k+1}$ 同痕的消失路径；其形成方式与\secref{16h}相同。设 $\widetilde\Delta_k,\widetilde\Delta_{k+1}$ 为与 $\gamma_k,\gamma_{k+1}$ 相关的 (B) 型子流形。
\StatementNumber{4}\begin{lemma}\label{st:18.4}
作为带标架 Lagrangian 球面，$\widetilde\Delta_{k+1}$ 与 $\tau_{\widetilde\Sigma_{\mu,-}}\tau_{\widetilde\Sigma_{\mu,+}}(\widetilde\Delta_k)$ 同痕。
\end{lemma}
\SourcePage{281}
\SourceFigure{figure-18-1}{0.469384110}{18.1}
\begin{proof}
以 $\widetilde\mu=\widetilde\mu_+\cup\widetilde\mu_-$ 及 $\widetilde\gamma_k,\widetilde\gamma_{k+1},\widetilde\gamma'_{k+1}$ 表示相应路径在 $\widetilde S$ 中的原像。它们都是 $\widetilde\pi$ 的匹配路径，而且
\[
\widetilde\gamma_{k+1}\simeq\widetilde\gamma'_{k+1}=\widetilde\mu_-\#(\widetilde\mu_+\#\widetilde\gamma_k).
\]
两次应用\stref{引理}{16.13}，即可得到所需结论，唯一可能剩下的是标架问题；但这也容易解决：对原 Lefschetz 纤维化 $E$ 的辛联络作变形，使 $\mu$ 的匹配圈可以用朴素方式构造，即同痕 $I$ 为常值。同样，对于匹配路径 $\widetilde\gamma_k,\widetilde\gamma_{k+1}$，同痕也为常值，这与证明\stref{引理}{16.13}时使用的 Lagrangian 连通和论证相容。
\end{proof}

\StatementNumber{5}\begin{note}\label{st:18.5}
以上三个引理都有将对合 $\iota_{\widetilde E}$ 考虑在内的加强版本：
\begin{itemize}
\item \stref{引理}{18.2}中的同痕自然等变，因为在局部模型中成立的是等式，而用于归约到该模型情形的变形与 $\iota_{\Efamily}$ 相容。
\item 对于\stref{引理}{18.3}，只须使用 Picard–Lefschetz 公式的等变版本。这里指出一个有用的附加细节：将该同痕与 $\tau_{\widetilde\Delta_k}|\widetilde\Delta_{k+1}$ 复合，得到等变微分同胚
\begin{equation}\label{eq:18.14}
\widetilde\Delta_{k+1}\longrightarrow\widetilde\Pi_k,
\end{equation}
因此，可以比较两个球面内不动点集的余定向。这些余定向自然对应于 $\widetilde S$ 中路径 $\widetilde\delta_k,\widetilde\gamma_{k+1}$ 的定向。结论是，两个余定向与\eref{18.14}相容，当且仅当它们在唯一公共点 $y=0$ 处的相交数 $\widetilde\delta_k\cdot\widetilde\gamma_{k+1}$ 为 $-1$。为看出这一点，取\stref{引理}{18.3}证明中的环路 $\lambda$，比较相应的 $\widetilde E$ 上的单值变换映射与
\SourcePage{282}
$\widetilde S$ 上的辫单值变换映射。一般而言，辫单值变换是保持 $\operatorname{Critv}(\widetilde\pi)$ 的 $\widetilde S$ 的微分同胚；在当前情形中，它就是沿 $\widetilde\gamma_k$ 的半扭转。若为 $\widetilde\gamma_{k+1}$ 选定定向，并为其像 $\widetilde\delta_k=t_{\widetilde\gamma_k}(\widetilde\gamma_{k+1})$ 赋予诱导定向，则局部相交数确实为 $-1$，见\sr{fig:18.2}{图~18.2}。
\SourceFigure{figure-18-2}{0.365538068}{18.2}
\item 对这一目的而言，我们对\stref{引理}{18.4}的证明不甚理想，因为每次只作一个 Dehn 扭转破坏了对称性。不过，适当改变论证便可轻易补救：先在 $E$ 中相对于 $M$ 构造 Lagrangian 球体的同痕 $\Delta_{k+1}\simeq\tau_{\Sigma_\mu}(\Delta_k)$，再将其提升到二重分歧覆盖。
\end{itemize}
\end{note}
